\section{Discussion}

\subsection{Why Production-Grounded Evaluation Matters}

Our results demonstrate that serial dependency introduces an evaluation axis invisible to static benchmarks.
When tasks are independent, an agent's failure at one task does not affect others, and the benchmark can measure each capability in isolation.
When tasks are serially dependent, failure propagates, and the evaluation framework must distinguish between ``cannot reach'' and ``cannot solve''---a distinction that resurrection enables.

This is not merely a technical detail but a fundamental shift in what evaluation frameworks measure: from \emph{can the agent solve task~$k$?} to \emph{can the agent progress along the dependency chain, and if not, where and why?}
Our finding that only 3 out of 22 configurations achieve zero-shot pipeline success---despite most models scoring highly on individual tasks---demonstrates that static benchmarks systematically overestimate production readiness.

\subsection{Resurrection as Methodology, Not Assistance}

A natural concern is that resurrection ``helps'' the agent by providing correct intermediate artifacts, thereby inflating scores.
We emphasize that resurrection is an \emph{evaluation protocol}, not a scoring adjustment.
The no-resurrection setting provides the ``strict'' measurement of zero-shot pipeline capability.
The resurrection setting provides the ``diagnostic'' measurement of node-level capability given correct prerequisites.
Both are reported independently; neither replaces the other.

\textbf{Potential biases.}
We acknowledge three potential biases: (1)~\emph{state discontinuity}---the golden artifact may differ in style from the agent's own output; (2)~\emph{selection bias}---only failed tasks are resurrected; (3)~\emph{moral hazard}---if agents were aware of resurrection, they might reduce early-stage effort.
We mitigate bias~(3) by ensuring agents are \emph{not informed} about resurrection---the mechanism is triggered externally by the evaluation orchestrator.

\subsection{Process Cost as Diagnostic Signal}

Our regression analysis reveals that process cost (build count, elapsed time) explains essentially zero percent of outcome variance ($R^2 = 0.007$).
This is not a failure of process metrics---it is a finding with direct implications for production deployment.
Agents that invest more computational budget do not systematically achieve better outcomes.
Two agents with identical pipeline scores may have radically different process profiles (efficient solver vs.\ persistent searcher), and this distinction is invisible to outcome-only evaluation but critical for deployment decisions regarding cost budgeting, timeout setting, and scaffold optimization.

\subsection{The Scaffolding Gap}

The performance delta between \yatcc{} and \yatcc{}-Hard quantifies the value of scaffolding in agentic coding tasks.
Top models on \yatcc{} (e.g., \texttt{kimi-k2.6} at PL=99.3) drop to PL=64.9 on \yatcc{}-Hard---a 34.4-point decline.
This reveals that scaffolded starter code provides substantial leverage, and that models appearing capable when filling in blanks may lack the architectural reasoning required for full from-scratch implementation.
The scaffolding gap thus serves as a diagnostic metric for distinguishing between \emph{completion capability} and \emph{construction capability}.

\subsection{Limitations}

\textbf{Domain specificity.}
\eva{} is grounded in compiler construction, a domain with particularly clean intermediate artifact specifications.
Other serial workflows---ML pipelines, DevOps chains, data engineering flows---may have less well-defined intermediate states, making resurrection harder to implement.
We do not claim that \eva{} generalizes to all serial settings, but argue that the \emph{evaluation methodology} (serial dependency + resurrection + process metrics) is transferable.

\textbf{Task count and statistical power.}
With six tasks, \eva{} provides limited statistical power for per-task analysis.
However, the serial structure means each task carries \emph{qualitatively different} information (different dependency depth, different failure propagation pattern), which compensates for the small task count with structural diversity.
Each model configuration is evaluated once per pipeline run; multi-seed evaluation is planned as future work.

\textbf{Contamination risk.}
\yatcc{} is a public GitHub repository, and frontier models may have been trained on its contents.
However, \eva{}'s key challenge is not reproducing known solutions but maintaining \emph{serial continuity}---producing correct intermediate artifacts that integrate across the pipeline.
Memorizing individual task solutions does not guarantee correct cross-task artifact flow.
Future versions could adopt temporal splitting or private codebases to mitigate this concern.

\subsection{Future Directions}


Multi-domain workloads: extending the evaluation methodology to ML engineering, DevOps, and data science pipelines.
Evolutionary gain measurement: comparing agent performance when inheriting self-produced vs.\ golden intermediate artifacts, quantifying how well agents ``understand their own code.''
Dynamic resurrection: allowing agents to request resurrection proactively, studying self-diagnosis capability.
Expanded model and scaffold coverage: evaluating additional frontier models and agent frameworks as they become available.
Infrastructure-aware workloads: incorporating deployment validation, runtime monitoring, and recovery scenarios.
